\section{Circuit Design}
\label{sec:circuit_design}

Present the transistor-level or block-level circuit implementation. Each
subsection should correspond to one major block in the architecture section.

\subsection{Main Circuit Block}

Describe the selected topology and explain why it satisfies the requirements
derived in Section~\ref{sec:design_considerations}. Include the schematic after
the design is finalized.

\begin{figure}[H]
    \centering
    \fbox{\parbox[c][5cm][c]{0.85\textwidth}{
        \centering
        Replace this box with the main circuit schematic.
    }}
    \caption{Schematic of the main circuit block.}
    \label{fig:main_circuit}
\end{figure}

\subsection{Biasing and Supporting Circuits}

Describe the bias circuit, common-mode generation, clock generation, reference
generation, or other supporting circuitry required by the design.

\subsection{Device Sizing}

Summarize the final device dimensions or component values. Add or remove columns
based on the project requirements.

\begin{table}[H]
    \centering
    \caption{Device sizing summary.}
    \label{tab:device_sizing}
    \renewcommand{\arraystretch}{1.2}
    \begin{tabular}{c c c c}
        \toprule
        \textbf{Device} & \textbf{Width / Fin} & \textbf{Length} & \textbf{Multiplier} \\
        \midrule
        TBD & TBD & TBD & TBD \\
        TBD & TBD & TBD & TBD \\
        TBD & TBD & TBD & TBD \\
        \bottomrule
    \end{tabular}
\end{table}

\subsection{Implementation Notes}

Record practical design decisions that are useful for review, such as layout
concerns, parasitic sensitivity, startup behavior, stability margin, or simulation
shortcuts.
